\documentclass[12pt]{article}
\usepackage[utf8]{inputenc}
\usepackage[T1]{fontenc}
\usepackage{hyperref}
\title{Lule Language Resources for K-12 Curriculum Development}
\author{Compiled by Claude Sonnet 5 (claude-sonnet-5) for Andr\'as Kornai}
\date{}
\begin{document}
\maketitle
\section{Introduction}
Lule (ISO 639-3 ule) is an extinct language of northwestern Argentina (Salta, Tucum\'an, Santiago del Estero). Its classification is disputed: some link it to Vilela, while others (Zamponi 2008) treat both as isolates. Whether Lule and Tonocot\'e were one language is also debated. Only a small body of data survives.
\section{Primary Educational Resources}
No school materials exist. No teaching resources were found.
\section{Linguistic and Reference Materials}
The main source is Antonio Machoni's \emph{Arte y vocabulario de la lengua lule y tonocot\'e} (Madrid, 1732), a missionary grammar and vocabulary; I did not locate a full-text copy. Overview: \url{https://en.wikipedia.org/wiki/Lule_language}; related: \url{https://en.wikipedia.org/wiki/Tonocot%C3%A9_language}. Scholarly pages: \url{https://uba.academia.edu/PedroViegasBarros} and \url{https://periodicos.sbu.unicamp.br/ojs/index.php/liames/article/download/1400/1390/1498}. I did not verify the contents of these two.
\section{Religious and Community Resources}
Machoni's work was written for evangelization. No community organizations were found.
\section{Digital and Audio Resources}
None found. There are no known recordings.
\section{Curriculum Development Context}
With no speakers and only a colonial-era grammar and word list, any heritage curriculum would focus on history, place names and the people's story rather than language instruction. Reconstruction claims would need caution.
\end{document}